\documentclass[11pt]{article}
\usepackage[margin=1in]{geometry}
\usepackage{amsmath,amssymb,amsthm}
\usepackage[round]{natbib}
\usepackage{hyperref}
\usepackage{microtype}

\newtheorem{lemma}{Lemma}
\newtheorem{remark}{Remark}

\title{Conformal Prediction and the Finite (Signed) de Finetti Representation}
\author{Peter Cotton}
\date{Draft --- companion note}

\begin{document}
\maketitle

\begin{abstract}
Conformal prediction needs only exchangeability, and exchangeability has a representation
theorem: de Finetti's. This note connects the two through the \emph{finite} form, in which
the mixing measure is allowed to be \emph{signed} \citep{kerns2006}. The link is modest but,
as far as a literature search reveals, unstated: the sign of the de Finetti mixing measure is
the sign of the term that controls whether split conformal's coverage is conditionally
adaptive or conditionally perverse, while the marginal guarantee is untouched throughout.
Proper (positive, extendable) mixtures make conformal conditionally self-correcting; the
signed, non-extendable corner---exactly the regime of constrained scores (de-meaned, ranked,
compositional)---makes it anti-adaptive. None of this changes the marginal coverage number;
it changes how much that number hides. I claim no new theorem of consequence: the
marginal-invariance is classical, the independence case of the lemma below is the known Beta
law for calibration-conditional coverage, and the dependence calculus is standard association
theory. What is offered is the bridge---reading conformal's conditional behaviour through the
sign of the finite de Finetti representation.
\end{abstract}

\section{Setup}

Split conformal. Scores $S_1,\dots,S_n$ (calibration) and $S_{n+1}$ (test) are jointly
exchangeable with a continuous joint law (no ties). The conformal width is the order
statistic $\hat q = S_{(k)}$ of the calibration scores, $k=\lceil(n+1)(1-\alpha)\rceil$, and a
test point is covered when $S_{n+1}\le\hat q$.

Two kinds of dependence live in any such problem, and they should not be confused.
\emph{Within a row}: inside a single $(x,y)$, does the residual's shape depend on $x$? This is
the subject of the companion paper \citep{cotton_marginally}, where the log-score regret of a
single-shape conformal system equals the mutual information $I(R;X)$; it is a property of one
observation's law and is present even when the rows are i.i.d. \emph{Across rows}: how do
distinct observations relate to one another---independent, sharing a latent cause, or mutually
constrained? This is the subject of de Finetti's theorem, and of this note. The two axes are
orthogonal; this note is entirely about the second.

\section{Finite de Finetti and the sign of the mixture}

de Finetti's theorem (infinite form) says an infinitely exchangeable sequence is a
\emph{positive} mixture of i.i.d.\ sequences: nature draws a parameter from a prior, then
generates i.i.d.\ data. A one-line variance identity shows the cost of that ``positive'': a
positive mixture has $\operatorname{Cov}(S_i,S_j)=\operatorname{Var}\big(\mathbb
E[S_i\mid\Theta]\big)\ge 0$, so proper de Finetti mixtures force non-negative correlation.
Finite exchangeable sequences are not so constrained: exchangeability only forces
$\rho\ge -1/(n-1)$ \citep{aldous1985}, and the negative band is real---sampling without
replacement sits at its floor.

The finite theory closes this gap in two ways. \citet{diaconis1980} show a finite exchangeable
sequence is \emph{approximately} a positive i.i.d.\ mixture, within total variation $O(k/n)$.
\citet{kerns2006} make it \emph{exact} at the price of signing: every finite exchangeable
sequence is a mixture of i.i.d.\ product measures with a \emph{signed} mixing measure (see
also \citealp{leonetti2018} for the partially-exchangeable extension). Geometrically---this
phrasing is the author's synthesis, not a quoted result---the i.i.d.\ product laws trace a
curve; positive mixtures fill its convex hull (the extendable, $\rho\ge 0$ laws); signed
mixtures fill its affine span (everything). A signed representation is thus the signature of
\emph{non-extendability}: a finite exchangeable sequence that cannot be continued to an
infinite one, and it is exactly the negative-correlation regime.

So the sign of the de Finetti mixing measure is a clean label: \textbf{positive} = extendable,
non-negatively correlated, a genuine prior; \textbf{signed} = finite, negatively correlated,
no genuine prior.

\section{Marginal coverage ignores the sign}

For continuous exchangeable scores the rank of $S_{n+1}$ among all $n+1$ is uniform, and
$S_{n+1}\le S_{(k)}$ iff that rank is at most $k$, so
\begin{equation}
\mathbb P\big(S_{n+1}\le\hat q\big)=\frac{k}{n+1}\ge 1-\alpha .
\end{equation}
This uses nothing but exchangeability \citep{vovk2005}. Positive or signed, the marginal
number is the same. The sign enters only when we ask what the certificate hides.

\section{The conditional-coverage decomposition}

Condition on the realised width. Write $H(s,q)=\mathbb P(S_{n+1}\le s\mid\hat q=q)$, so coverage
conditional on the width is the diagonal $c(q)=H(q,q)$.

\begin{lemma}[Conditional-coverage slope]\label{lem:slope}
If $(S_{n+1},\hat q)$ has a joint density, then
\begin{equation}
c'(q)=\underbrace{f_{S_{n+1}\mid\hat q}(q\mid q)}_{\textup{(i) threshold}}
\;+\;\underbrace{\partial_q\,\mathbb P(S_{n+1}\le q\mid\hat q=q)}_{\textup{(ii) dependence}} .
\end{equation}
Term \textup{(i)} is non-negative. Term \textup{(ii)} is $\le 0$ when $S_{n+1}$ is
stochastically increasing in $\hat q$, $=0$ under independence, and $\ge 0$ when stochastically
decreasing.
\end{lemma}

\begin{proof}
Total derivative of $H(q,q)$ along the diagonal: $c'(q)=\partial_s H(s,q)\big|_{s=q}+\partial_q
H(s,q)\big|_{s=q}$. The first equals $f_{S_{n+1}\mid\hat q}(q\mid q)\ge 0$. For the second,
$\partial_q H(s,q)\le 0$ for all $s$ exactly when $\mathbb P(S_{n+1}\le s\mid\hat q=q)$ is
non-increasing in $q$, the definition of $S_{n+1}$ being stochastically increasing in
$\hat q$; the reverse inequality is stochastic decrease.
\end{proof}

\begin{remark}[The independence case is classical]
Under independence term (ii) vanishes and $c(q)=F(q)$, the test-score CDF, which already rises
with the width purely from estimation. This is the known law of calibration-conditional
coverage: $F(\hat q)=F(S_{(k)})\sim\operatorname{Beta}(k,n-k+1)$
\citep{vovk2012,marques2024,angelopoulos2023gentle}. Lemma~\ref{lem:slope} is its extension to
dependent scores, with term (ii) the new piece.
\end{remark}

\begin{remark}[Reading the sign]
The width $\hat q=S_{(k)}$ is a coordinatewise non-decreasing function of the calibration
scores, and the test score is disjoint from them. A positive (proper) de Finetti mixture makes
the score vector \emph{associated} \citep{esary1967}; association is preserved by monotone
maps, so $(\hat q,S_{n+1})$ are positively dependent, term (ii) is $\le 0$, and it cancels
against the threshold term---conformal is \emph{conditionally adaptive}, flatter than
independence, because pooling the calibration residuals silently conditions on the realised
world. A signed, negatively-correlated law makes the vector \emph{negatively associated}
\citep{joagdev1983}---the property of without-replacement samples, multinomials, permutation
and rank laws, and Dirichlet/compositional vectors, preserved under monotone maps of disjoint
subsets---so term (ii) is $\ge 0$ and adds to the threshold term: conformal is
\emph{conditionally anti-adaptive}. The calibration set anti-predicts the test, so coverage
over-covers when the calibration scores were large and under-covers when small, even as the
marginal stays pinned. The discrete (order-statistic) version replaces the derivative by a
difference and behaves identically; the density assumption is for readability.
\end{remark}

\section{Why the signed corner is not exotic: constrained scores}

The negative corner arises whenever the nonconformity score is defined \emph{relative to the
batch}, because a linear constraint forces it. De-meaning ($\sum_i s_i=0$), differencing,
ranks or percentiles-within-batch, and compositional or budgeted targets all impose such a
constraint, and the constraint \emph{is} the negative association \citep{joagdev1983}. For
these scores conformal remains marginally valid---if the relativisation is symmetric across
calibration and test, exchangeability is preserved---but its conditional coverage is the
fanning, anti-adaptive kind of Lemma~\ref{lem:slope}.

A caution: if the relativisation treats the test point asymmetrically (de-meaning by a
training-only mean, say), exchangeability itself fails and even the marginal guarantee is no
longer assured. This is the familiar reason split conformal uses out-of-fold rather than
in-sample residuals \citep{lei2018}; in-sample regression residuals are both constrained
($\sum_i e_i=0$) and heteroscedastic (variance $1-h_{ii}$), and so are not exchangeable to
begin with.

\section{Scope, and what is and is not new}

The contribution is a viewpoint, not a theorem. To be explicit:
\begin{itemize}
\item \emph{Classical, cited here:} marginal coverage under exchangeability
\citep{vovk2005}; the Beta law of calibration-conditional coverage
\citep{vovk2012,marques2024}; training-conditional coverage as a studied object
\citep{bian2023}; coverage degradation under non-exchangeability bounded by
total-variation/mixing magnitude \citep{barber2023beyond,oliveira2024}; positive and negative
association \citep{esary1967,joagdev1983}; finite and signed de Finetti
\citep{diaconis1980,kerns2006,leonetti2018}.
\item \emph{Closest prior work:} \citet{barberpananjady2026} show conformal can both under- and
over-cover under temporal dependence. That analysis is marginal and indexes departures by
dependence \emph{magnitude} (a switch coefficient / $\beta$-mixing), not by \emph{sign}, and
does not treat calibration-conditional behaviour. The present note is complementary: the sign,
read through the de Finetti representation, and the conditional rather than marginal object.
\item \emph{Offered here:} the identification of the conditional-coverage dependence term's
sign with the sign of the finite de Finetti mixing measure, and the consequent reading of
constrained/relative scores (the signed corner) as the regime where conformal's marginal
certificate is least informative about conditional reliability.
\end{itemize}

The link runs only one way and only so far. de Finetti is a representation of a law; conformal
is a procedure run without knowing the law. A signed representation does not yield a usable
predictor---to use any de Finetti measure one must estimate it, which is modelling, with no
distribution-free finite-sample guarantee. So the note does not claim de Finetti
``supersedes'' or ``contains'' conformal; recent work even argues conformal is not a de
Finetti/Bayesian conditional at all \citep{datta2025}. The representation diagnoses conformal's
conditional behaviour; it does not replace it.

\begin{thebibliography}{99}
\bibitem[Aldous(1985)]{aldous1985} Aldous, D.~J. (1985). Exchangeability and related topics.
\emph{\'Ecole d'\'Et\'e de Probabilit\'es de Saint-Flour XIII}, Lecture Notes in Mathematics
1117, Springer.
\bibitem[Angelopoulos and Bates(2023)]{angelopoulos2023gentle} Angelopoulos, A.~N. and Bates,
S. (2023). A gentle introduction to conformal prediction and distribution-free uncertainty
quantification. \emph{Foundations and Trends in Machine Learning}. arXiv:2107.07511.
\bibitem[Barber et~al.(2023)]{barber2023beyond} Barber, R.~F., Cand\`es, E.~J., Ramdas, A. and
Tibshirani, R.~J. (2023). Conformal prediction beyond exchangeability. \emph{The Annals of
Statistics} 51(2):816--845.
\bibitem[Barber and Pananjady(2026)]{barberpananjady2026} Barber, R.~F. and Pananjady, A.
(2026). Predictive inference for time series: why is split conformal effective despite temporal
dependence? Preprint, arXiv:2510.02471.
\bibitem[Bian and Barber(2023)]{bian2023} Bian, M. and Barber, R.~F. (2023).
Training-conditional coverage for distribution-free predictive inference. \emph{Electronic
Journal of Statistics} 17(2):2044--2066.
\bibitem[Cotton()]{cotton_marginally} Cotton, P. Marginally Useful: Formalizing the Information
Gap in Conformal Prediction. Companion paper.
\bibitem[Datta et~al.(2025)]{datta2025} Datta, J., Polson, N.~G., Sokolov, V. and Zantedeschi,
D. (2025). Conformal Prediction = Bayes? Preprint, arXiv:2512.23308.
\bibitem[Diaconis and Freedman(1980)]{diaconis1980} Diaconis, P. and Freedman, D. (1980).
Finite exchangeable sequences. \emph{The Annals of Probability} 8(4):745--764.
\bibitem[Esary et~al.(1967)]{esary1967} Esary, J.~D., Proschan, F. and Walkup, D.~W. (1967).
Association of random variables, with applications. \emph{The Annals of Mathematical
Statistics} 38(5):1466--1474.
\bibitem[Joag-Dev and Proschan(1983)]{joagdev1983} Joag-Dev, K. and Proschan, F. (1983).
Negative association of random variables, with applications. \emph{The Annals of Statistics}
11(1):286--295.
\bibitem[Kerns and Sz\'ekely(2006)]{kerns2006} Kerns, G.~J. and Sz\'ekely, G.~J. (2006). De
Finetti's theorem for abstract finite exchangeable sequences. \emph{Journal of Theoretical
Probability} 19(3):589--608.
\bibitem[Leonetti(2018)]{leonetti2018} Leonetti, P. (2018). Finite partially exchangeable laws
are signed mixtures of product laws. \emph{Sankhy\=a A} 80(2):195--214.
\bibitem[Lei et~al.(2018)]{lei2018} Lei, J., G'Sell, M., Rinaldo, A., Tibshirani, R.~J. and
Wasserman, L. (2018). Distribution-free predictive inference for regression. \emph{Journal of
the American Statistical Association} 113(523):1094--1111.
\bibitem[Marques F.(2024)]{marques2024} Marques F., P.~C. (2024). Universal distribution of the
empirical coverage in split conformal prediction. Preprint, arXiv:2303.02770.
\bibitem[Oliveira et~al.(2024)]{oliveira2024} Oliveira, R.~I., Orenstein, P., Ramos, T. and
Romano, J.~V. (2024). Split conformal prediction and non-exchangeable data. \emph{Journal of
Machine Learning Research} 25.
\bibitem[Vovk(2012)]{vovk2012} Vovk, V. (2012). Conditional validity of inductive conformal
predictors. \emph{Asian Conference on Machine Learning}, PMLR 25:475--490.
\bibitem[Vovk et~al.(2005)]{vovk2005} Vovk, V., Gammerman, A. and Shafer, G. (2005).
\emph{Algorithmic Learning in a Random World}. Springer.
\end{thebibliography}

\end{document}
